% P11 paper module: R29 small-source constraint-normal mismatch.

\subsection{Small-source constraint-normal mismatch}
\label{subsec:o3ab-small-source-normal}

The fixed R28 normal obstruction can be decided on a nontrivial concrete region of source triples.  Keep $0<R<S<T_0$ and the notation of Subsection~\ref{subsec:o3aa-fixed-gamma-crossblock}.  Thus $r_X$ is the Riesz vector of the transported first boundary functional in the $T_0$-whitened source Hilbert space, and
\[
 s_{R,S,T_0}:=(I-WW^*)r_S.
\]

\begin{proposition}[Dual norm of the first-boundary normal]
\label{prop:o3ab-dual-normal}
For $X\in\{R,S\}$,
\[
 \boxed{
 \|r_X\|
 =\sup_{0\ne f\in\KXR{X}^-}
 \frac{|\beta_X^{(0)}(f)|}
 {q_{T_0}^X(J_{X,T_0}f)^{1/2}}.
 }
\tag{O3AB.1}\label{eq:o3ab-dual-normal}
\]
Moreover
\[
 \boxed{
 \|r_R\|^2\le\frac23R^3.
 }
\tag{O3AB.2}\label{eq:o3ab-small-R-bound}
\]
\end{proposition}

\begin{proof}
If $x=B_X^{1/2}f$, then
\[
 \|x\|^2=\langle B_Xf,f\rangle
 =q_{T_0}^X(J_{X,T_0}f),
\]
and the defining Riesz identity for $r_X$ gives
\[
 \langle x,r_X\rangle=\beta_X^{(0)}(f).
\]
Taking the dual norm proves~\eqref{eq:o3ab-dual-normal}.

Now
\[
 \beta_R^{(0)}(f)
 =\int_{-R}^R\phi_R(u)f(u)\,du,
 \qquad
 \phi_R(u):=\operatorname{sgn}(u)I_0(|u|),
\]
with
\[
 I_0(r)=2(1-e^{-r/2})\le r.
\]
By positivity of the Schur term, Gamma compatibility, and
$C_{\Gamma,R}\ge I$,
\[
 q_{T_0}^X(J_{R,T_0}f)
 \ge\mathfrak c_{\Gamma,R}[f]
 \ge\|f\|_2^2.
\]
Hence
\[
 \|r_R\|\le\|\phi_R\|_2
 \le\left(2\int_0^Rr^2dr\right)^{1/2},
\]
which is~\eqref{eq:o3ab-small-R-bound}.
\end{proof}

For reference, the elementary exact kernel norm is
\[
 \|\phi_R\|_2^2
 =8\left[R-4(1-e^{-R/2})+(1-e^{-R})\right]
 =\frac23R^3+O(R^4).
\tag{O3AB.3}\label{eq:o3ab-exact-kernel-norm}
\]

\begin{theorem}[Concrete small-source normal mismatch]
\label{thm:o3ab-small-source-mismatch}
Fix $0<S<T_0$.  Then $\|r_S\|>0$, and for every
\[
 0<R<\min\left\{S,
 \left(\frac32\|r_S\|^2\right)^{1/3}\right\}
\]
one has
\[
 \boxed{s_{R,S,T_0}\ne0.}
\tag{O3AB.4}\label{eq:o3ab-s-nonzero}
\]
More quantitatively,
\[
 \boxed{
 \|s_{R,S,T_0}\|^2
 \ge\|r_S\|^2-\frac23R^3.
 }
\tag{O3AB.5}\label{eq:o3ab-s-lower}
\]
\end{theorem}

\begin{proof}
The first boundary functional is nonzero on the odd source space: choose a smooth odd $h$ with $h\ge0$ and $h\not\equiv0$ on the positive half-interval.  Then
\[
 \beta_S^{(0)}(h)
 =2\int_0^SI_0(r)h(r)dr>0,
\]
so its Riesz vector $r_S$ is nonzero.

By~\eqref{eq:o3aa-normal-norm},
\[
 \|s_{R,S,T_0}\|^2
 =\|r_S\|^2-\|r_R\|^2.
\]
Use~\eqref{eq:o3ab-small-R-bound} to obtain
\eqref{eq:o3ab-s-lower}.  Positivity of the right-hand side gives
\eqref{eq:o3ab-s-nonzero}.
\end{proof}

\begin{corollary}[Persistent full inverse-root defect]
\label{cor:o3ab-inverse-defect}
In the region of Theorem~\ref{thm:o3ab-small-source-mismatch},
\[
 \boxed{
 \|D_\infty^-(R,S,T_0)\|
 \ge
 \sqrt{1-\frac{2R^3}{3\|r_S\|^2}}
 >0.
 }
\tag{O3AB.6}\label{eq:o3ab-D-lower}
\]
For fixed $S<T_0$,
\[
 \boxed{
 \liminf_{R\downarrow0}
 \|D_\infty^-(R,S,T_0)\|\ge1.
 }
\tag{O3AB.7}\label{eq:o3ab-D-small-R}
\]
\end{corollary}

\begin{proof}
Combine Theorem~\ref{thm:o3ab-small-source-mismatch} with the R28 estimate
\eqref{eq:o3aa-normal-lower}.  Since
$\|r_R\|\to0$ as $R\downarrow0$, the orthogonal identity
$\|s\|^2=\|r_S\|^2-\|r_R\|^2$ gives
$\|s\|/\|r_S\|\to1$.
\end{proof}

\begin{remark}[Computable sufficient threshold]
\label{rem:o3ab-computable-threshold}
For any fixed smooth odd $h\in C_c^\infty((-S,S))$ with
$\beta_S^{(0)}(h)\ne0$, Proposition~\ref{prop:o3ab-dual-normal} gives
\[
 c_h:=\frac{|\beta_S^{(0)}(h)|}
 {q_{T_0}^X(J_{S,T_0}h)^{1/2}}
 \le\|r_S\|.
\]
Hence the entirely fixed-window condition
\[
 R<\min\left\{S,\left(\frac32c_h^2\right)^{1/3}\right\}
\]
already implies $s_{R,S,T_0}\ne0$.
\end{remark}

\begin{remark}[R29 firewall]
\label{rem:o3ab-firewall}
This is a concrete P11 no-go for full inverse-root intertwining on a whole small-source region; it is not an abstract countermodel.  Nevertheless it remains on the modulus/inverse-functional-calculus layer.  By the R14 polar-gauge firewall, nonvanishing of $D_\infty^-$ does not by itself imply nonconvergence of the actual normalized future transport, nonvanishing of the R22 angle defect, or failure of a global Object~X construction.
\end{remark}

\begin{openproblem}[R29-F: full fixed-normal classification]
\label{open:o3ab-normal-classification}
Determine whether
\[
 s_{R,S,T_0}\ne0
\]
for every strict inclusion $0<R<S<T_0$, or characterize exactly those fixed triples for which the baseline Riesz normal of the first boundary functional lies entirely in the normalized old-source range.  The tangential rank-one test of R28 remains a separate question.
\end{openproblem}
